\section{Stack tecnológico seleccionado}

El stack tecnológico es el conjunto de herramientas, programas, aplicaciones u otras piezas de software que son utilizadas para la creación de un producto tecnológico. Una vez estudiadas las principales opciones para el desempeño de cada sector del proyecto, se cuenta con una visión global lo suficientemente rica como para hacer una elección que se ajuste a las necesidades y preferencias del proyecto y su desarrollador.

En esta sección se expone cada punto de esta selección:

\begin{itemize}
    \item Para el desarrollo del frontend de esta aplicación se ha optado por \textbf{Next.js} debido a sus ventajas en rendimiento, optimización para SEO y compatibilidad con React, que es una de las tecnologías más populares e innovadoras hoy en día, por lo que dominarlo resulta de gran interés. Entre los aspectos más atractivos de esta opción se encuentra también que cuenta con una plataforma de lanzamiento y hosting de aplicaciones gratuito, rápido y simple: Vercel; además de que la configuración inicial de un proyecto web se hace muy sencilla y existe una documentación rica, una comunidad activa y numerosos recursos para facilitar el desarrollo web.
    \item En el backend, se ha elegido \textbf{Node.js con Express} por su flexibilidad y el beneficio de utilizar el mismo lenguaje de programación, JavaScript, en todo el stack de desarrollo. Además, se trata de una tecnología ampliamente utilizada en el desarrollo de microservicios y cuenta con librerías que facilitan el acceso al tipo de base de datos que se utilizará.
    \item En cuanto a la base de datos, se ha decidido emplear \textbf{MongoDB} por su modelo flexible y su capacidad para manejar estructuras de datos dinámicas y en formato JSON, el mismo en que se manejan las estructuras de información en JavaScript. También, al ser las bases de datos relacionales las más trabajadas durante el grado, resulta de especial interés cultivar nuevas habilidades con modelos no relacionales como el de MongoDB.
    \item Para el consumo de APIs se usará \textbf{Axios} frente a FetchAPI por simplicidad y aceleración del desarrollo, dado que Axios ofrece más funcionalidades listas para usar como la conversión automática a formato JSON, el manejo de tokens o el manejo de errores integrado; mientras que Fetch, aunque es más ligero, requiere más código manual para tareas frecuentes. Dado que en esta aplicación se implementará autentifación de usuarios y una cantidad considerable de funciones, Axios se presenta como la opción más cómoda y compatible con el resto del stack.
    \item Como herramienta de pruebas de la API, se utiliza \textbf{Postman} frente a Insomnia por una razón principal: la comodidad de tenerlo ya instalado y estar familiarizada, al estar trabajando con ello en otros proyectos personales.
\end{itemize}
\begin{comment}
    \item Las inteligencias artificiles generativas que se han empleado durante la construcción del proyecto son ChatGPT para temas de investigación y clarificación de documentación, y el chat integrado de Github Copilot en Visual Studio Code por la comodidad que ofrece al trabajar directamente sobre el contexto completo del código, haciendo más eficiente y completo tanto el debugging como la entrega de propuestas de optimización y otras mejoras de los componentes y páginas.
    Por el momento, no se ha escogido ninguna plataforma de computación en la nube porque el alcance del proyecto no requiere todavía un despliegue en producción que incluya el uso de servicios externos.
\end{comment}


\subsubsection{Calidad del Código}
Adicionalmente, tanto en el frontend como en el backend, se han empleado varias herramientas destinadas a asegurar la calidad del código:

\begin{itemize}
    \item ESLint: herramienta de análisis estático configurada con reglas para mantener consistencia de código, es decir, un estilo limpio y homogéneo entre los distintos componentes, páginas y otros archivos auxiliares. También ayuda a detectar posibles errores en el código, como variables o importaciones no utilizadas.\cite{eslint_docs}
    \item TypeScript: proporciona verificación de tipos de datos en tiempo de compilación, lo cual permite detectar incompatibilidades en las propiedades de los componentes, llamadas a funciones mal tipadas, o en ocasiones el uso incorrecto de estados.
    \item Prettier: se trata de un formateador automático de código que garazantiza una aparencia consistende de este. Se han configurado reglas para limitar la longitud de línea (\texttt{printWidth: 80}), definir una indentación consistente (\texttt{tabWidth: 2}), forzar el punto y coma al final de las sentencias (\texttt{semi: true}), y estandarizar el uso de comillas dobles (\texttt{singleQuote: false}), entre otras.\cite{prettier_docs}
\end{itemize}